\documentclass{standalone}
\usepackage[european]{circuitikz}
\usepackage[scaled]{helvet}
\usepackage[T1]{fontenc}
\renewcommand\familydefault{\sfdefault}


\begin{document}

\begin{circuitikz}[anchor=center]

\ctikzset{bipoles/generic/height=0.1}
\ctikzset{bipoles/generic/width=0.3}
    
\draw (0,0) node[dipchip,
    num pins=8,
    external pins width=0.9,
    scale = 2,
    external pad fraction=4,
    align=center,](C){};
    
\draw (C.west) node[anchor=north,rotate=90, inner sep=1.5em]{\Large UASTW};
\draw (C.east) node[anchor=north,rotate=-90, inner sep=1.5em]{\Large DPS368};
    
\draw (C.west) node[anchor=north,rotate=90, inner sep=1.5em]{\Large UASTW};
\draw (C.east) node[anchor=north,rotate=-90, inner sep=1.5em]{\Large DPS368};
    
\draw[-latex] (C.pin 1) -- ++(-1,0) node[anchor=east]{GND};
\node [left,anchor=east] at (C.bpin 1) {GND};

\draw[dashed] (C.pin 2) to[short,-*] ++(-0.5,0) node[inner sep=0](intcon) {};
\draw[dotted] (intcon) to[R, european resistor] node[pos=0.5,anchor=east]{R\textsubscript{3}} ++(0,-0.7) node[scale=0.4,ground]{};
\draw[-latex,dashed] (intcon) -- ++(-0.5,0) node[anchor=east](SDO){INT};
\node [left,anchor=east] at (C.bpin 2) {SDO/INT/AD0};

\draw[latex-] (C.pin 3) -- ++(-1,0) node[anchor=east] {SCL};
\node [left,anchor=east] at (C.bpin 3) {SCL/SCK};

\draw[latex-] (C.pin 4) -- ++(-1,0) node[anchor=east] {SDA};
\node [left,anchor=east] at (C.bpin 4) {SDI/SDA};

\draw (C.pin 5) -| ++(1,-1) node[ground] {};
\node [right,anchor=west] at (C.bpin 5) {GND};

\draw[latex-] (C.pin 7) -- ++(1,0) node[anchor=west] {VDDIO};
\node [right,anchor=west] at (C.bpin 7) {VDDIO};

\draw[latex-] (C.pin 6) -- ++(1,0) node[anchor=west] {Not Connected};
\node [right,anchor=west] at (C.bpin 6) {CSB};



\draw[latex-] (C.pin 8) -- ++(1,0) node[anchor=west] {VDD};
\node [right,anchor=west] at (C.bpin 8) {VDD};

\draw(C.south west) ++(0,-1.2) node[anchor=north west, align=left, text width=9cm]{The pulldown resistor R\textsubscript{3} (100 k$\Omega$) is optional and will set the address to 0x76 instead of 0x77.};


\end{circuitikz}

\end{document}
